\documentclass[10pt,a4paper]{article}
\usepackage[margin=2.1cm]{geometry}
\usepackage{amsmath,amssymb}
\usepackage{booktabs,tabularx,array}
\usepackage[T1]{fontenc}
\usepackage{lmodern}
\usepackage{microtype}
\usepackage{xcolor}
\usepackage[colorlinks=true,linkcolor=blue!50!black,urlcolor=blue!50!black]{hyperref}
\usepackage{enumitem}
\setlist{nosep,leftmargin=1.4em}
\newcommand{\code}[1]{\texttt{\small #1}}
\newcommand{\dd}{\,\mathrm{d}}
\newcommand{\sinc}{\operatorname{sinc}}
\newcommand{\std}{\operatorname{std}}
\newcommand{\med}{\operatorname{median}}
\newcommand{\caveat}[1]{\par\smallskip\noindent\colorbox{orange!10}{\parbox{\dimexpr\linewidth-2\fboxsep}{\small\textbf{Assumption / status.} #1}}\par\smallskip}

\title{From the readout trace to the amplifier physics:\\ every formula, as implemented}
\author{ccd\_opt\_ml amplifier-physics study (PA03 boards 101/102, LSM ACM106/109/110)}
\date{17 November 2026; setup section added 20 November 2026}

\begin{document}
\maketitle

\noindent This note writes down, step by step, how each physics number in the two-page summary is computed from the raw
traces. Every equation is the one the code evaluates (script names in \code{monospace}, all in \code{amp\_phys/}).
Where the code makes an approximation or an assumption, a shaded box says so. Image-side inputs (the image gain
$g$ and image read noise $\sigma_e$ from the peak fits) are taken as given.

\tableofcontents

%=====================================================================
\section{Constants and notation}

\begin{center}\small
\begin{tabularx}{\linewidth}{@{}llX@{}}\toprule
symbol & value & meaning / source\\\midrule
$f_s$ & $15\times10^6$\,s$^{-1}$ & ADC sample rate (LTC2387) \\
$v[n]$ & ADU & raw sample $n$ of one channel \\
$b_k[n]$ & 0/1 & sequencer flag bit $k$: 0 pedestal window, 1 signal window, 4 low during reset, 6 H transfer, 7 V transfer\\
$\tau$ & 10\,\textmu s (150 samples) nominal & CDS integration window\\
$\Delta t$ & 13.8\,\textmu s (207 samples) & pedestal-to-signal window centre spacing\\
$T$ & 32.6\,\textmu s (489 samples) & skip period\\
$G$ & 0.875\,ADU/\textmu V & chain gain at the video input: $27.34\ \text{V/V} / 31.25\ \text{\textmu V/LSB}$ (design)\\
$g$ & ADU/e$^-$ & image gain (peak fit)\\
$A$ & 0.85 & source-follower gain (layout-based, model)\\
$C_{\rm ox}$ & 0.454\,fF/\textmu m$^2$ & ON stack: $\varepsilon_0/(t_{\rm ox}/3.9+t_{\rm n}/7.5)$, 50\,nm + 50\,nm\\
$t_{\rm eq}$ & 76\,nm & oxide-equivalent thickness $50+50\cdot3.9/7.5$\\
$WL$ & 345\,\textmu m$^2$ & output-FET gate area (layout, 330--360)\\
$W_{\rm RG}L_{\rm RG}$ & $6\times6$\,\textmu m$^2$ & reset-FET gate (layout)\\
$T_{\rm ccd}$ & 130\,K (PA03), 180\,K (LSM) & \\
$\lambda_t$ & 0.1\,nm & oxide tunnelling length (McWhorter)\\
$q,\ k_B$ & & elementary charge, Boltzmann constant\\\bottomrule
\end{tabularx}
\end{center}

\caveat{$G=0.875$ is the design value. The earlier dead-channel Johnson calibration (0.853) is withdrawn: board-102
ch0--2 have transistors and are not open inputs. Every conversion from ADU to volts carries this design-value
uncertainty (the ACM input buffer gain is still ``to be checked''). Quantities in electrons ($\div g$) or in
capacitance via $q/g$ do not depend on $G$ or $A$ at all.}

%=====================================================================
\section{From raw samples to skips (\code{trace\_anatomy.py})}

\paragraph{Cycle segmentation.} Rising edges of bit 0 give the pedestal-window starts $e_j$. The period is
$P=\med(e_{j+1}-e_j)$ samples. A cycle $j$ is a real skip if it is one period long ($|e_{j+1}-e_j-P|\le3$) and is
followed by a signal window (bit 1 set at $e_j+s_0+2$, with $s_0$ the median signal-window offset). Consecutive cycles
one period apart belong to the same pixel; a gap starts a new pixel. The first 20 skips of every pixel are dropped
(first-skip transient). This leaves cycles $\mathbf{c}_j = (v[e_j],\dots,v[e_j+P-1])$.

\paragraph{Template and residual.}
\begin{equation}
\bar{\mathbf c}[m]=\med_j\,\mathbf c_j[m],\qquad \mathbf r_j=\mathbf c_j-\bar{\mathbf c}.
\end{equation}
The template $\bar{\mathbf c}$ is the deterministic waveform (reset pulse, summing-well dip, window levels); the
residual $\mathbf r_j$ is the noise.

\paragraph{CDS value per skip.} With the pedestal window length $L_w$ (from the flags) and 5 samples trimmed at each
end,
\begin{equation}
p_j=\frac{1}{L_w-10}\sum_{m=5}^{L_w-6}c_j[m],\qquad
s_j=\frac{1}{L_w-10}\sum_{m=s_0+5}^{s_0+L_w-6}c_j[m],\qquad
x_j=s_j-p_j .
\end{equation}

\paragraph{Within-pixel differences.} Every noise statistic uses consecutive skips of the same pixel, so the pixel
charge cancels:
\begin{equation}
d_j=x_{j+1}-x_j\quad(\text{only for } j,j+1 \text{ in the same pixel}).
\end{equation}

%=====================================================================
\section{Time-domain noise quantities}

\subsection{Single-skip noise $\sigma_1$ and its input density}
\begin{equation}
\sigma_1=\frac{\std(d_j)}{\sqrt2}\ \ [\text{ADU}],\qquad
\sigma_{1,e}=\frac{\sigma_1}{g}\ \ [e^-].
\end{equation}
For white noise of one-sided density $e_n$ (V/$\sqrt{\text{Hz}}$) a boxcar of length $\tau$ has variance
$e_n^2/(2\tau)$, and the CDS difference of two independent boxcars has $e_n^2/\tau$. Hence
\begin{equation}
e_n=\frac{\sigma_1}{G}\sqrt{\tau}\quad[\text{\textmu V}\sqrt{\text s}]\ \to\ \text{nV}/\sqrt{\text{Hz}}\ (\times10^3).
\label{eq:en}
\end{equation}
\caveat{The code uses $\tau=L_w/f_s$ in eq.~\eqref{eq:en} although the average runs over $L_w-10$ samples; the
quoted density is therefore $\sqrt{150/140}-1\approx3.5\%$ high. Also, $\sigma_1$ contains every noise source (white,
1/f, setup); calling $e_n$ a ``white density'' is exact only if the noise is white (Sec.~\ref{sec:H}).}

Why consecutive skips are differenced, how $\sigma_1$ differs from the textbook $\sigma_{\rm CDS}$ obtained
from a PSD, and why it is blind to setup noise, is in Sec.~\ref{sec:sigma1}.

These give the table rows \emph{noise 3.7--4.4\,e} ($\sigma_{1,e}$) and \emph{18.7--21.1\,nV/$\sqrt{\rm Hz}$}
($e_n$).

\subsection{Other per-skip statistics}
\begin{align}
\rho_1 &= \operatorname{corr}(\tilde x_j,\tilde x_{j+1}),\quad \tilde x_j = x_j-\langle x\rangle_{\text{pixel}}
&&\text{(lag-1 autocorrelation: slow noise)}\\
\sigma_{\rm lvl} &= \std(\ell_{j+1}-\ell_j)/\sqrt2,\quad \ell_j=\tfrac12(p_j+s_j)
&&\text{(per-skip level scatter: reset kTC + electronics)}\\
\kappa &= \frac{\langle (d-\bar d)^4\rangle}{\langle (d-\bar d)^2\rangle^2}-3
&&\text{(excess kurtosis of } d_j\text{; telegraph noise} \Rightarrow \kappa>0)
\end{align}
(pandas' bias-corrected estimator). For a Gaussian $\kappa=0$; a two-level telegraph signal with comparable white noise
gives $\kappa$ of order 1 and more. ACM109 ch1 has 2.7; 101 ch2 is $\approx0$ in the median trace but large in its
episodes (90th-percentile over traces, \code{skip\_averaging.py}).

\subsection{Noise after averaging $N$ skips (\code{skip\_averaging.py})}
Per pixel the CDS series $x_j$ is cut into consecutive blocks of $N$, the block means $X_k^{(N)}$ are differenced
within the pixel, after removing the mean within-pixel profile (a sequence-locked drift common to all pixels):
\begin{equation}
\sigma_N=\frac{\std\big(X^{(N)}_{k+1}-X^{(N)}_k\big)}{\sqrt2}, \qquad N=1,4,16,\dots,512.
\end{equation}
White noise gives $\sigma_N=\sigma_1/\sqrt N$; 1/f, telegraph and setup wander make $\sigma_N$ flatten.

\subsection{Image cross-check}
\begin{equation}
\mathcal R_{\rm img}=\frac{\sigma_e^{\rm image}(N_{\rm img})}{\sigma_{1,e}/\sqrt{N_{\rm img}}},
\qquad N_{\rm img}=500\ (\text{PA03}),\ 1000\ (\text{LSM}).
\end{equation}
$\mathcal R_{\rm img}\approx1$ means the image reaches the white-noise expectation (PA03: 0.84--0.89);
$\mathcal R>1$ means extra slow noise or charge arriving during skipping (LSM: 1.2--2.0).

%=====================================================================
\section{Frequency-domain quantities: the setup number}

\paragraph{Residual spectrum.} The residuals of all skip cycles are concatenated, $r=(\mathbf r_1,\mathbf r_2,\dots)$,
and a Welch estimate is taken (Hann window, $2^{15}$ samples per segment):
\begin{equation}
\text{ASD}(f)=\frac{\sqrt{S_{rr}(f)}}{G}\ \ [\text{nV}/\sqrt{\text{Hz}}],\qquad
\text{ASD}_{[f_1,f_2]}=\med_{f_1\le f<f_2}\text{ASD}(f).
\end{equation}
The \textbf{setup number} is $\text{ASD}_{[200\,\text{Hz},1\,\text{kHz}]}$ (also reported: 1--10\,kHz, 10--100\,kHz,
0.1--1\,MHz, 1--7.5\,MHz).
\caveat{The concatenation joins skip cycles, but drops the pixel boundaries and H/V transfers, so low-frequency
content near $1/(\text{trace length})$ depends on the trace length. Comparing setups therefore needs equal trace
lengths: Oct.\ 0.2\,s traces 48.8 vs Nov.\ 34.8\,\textmu V/$\sqrt{\rm Hz}$ ($-29\%$); the 1.1\,s Oct.\ traces (68.2)
are not comparable.}
What the ASD is, why a clocked-trace ASD cannot be integrated into a $\sigma$, and what isolates a setup number from
the amplifier: Sec.~\ref{sec:setup2}.

\paragraph{Lines.} $E(f)=S_{rr}(f)/\widetilde S_{rr}(f)$ with $\widetilde S$ a 31-bin running median; the six largest
$E$ above 20\,kHz are listed (e.g.\ the 30.6\,kHz comb on ACM110 ch3).

\paragraph{Bursts, stationarity, drift.} The residual is cut into $\approx1$\,ms blocks $B_i$ (a whole number of
cycles):
\begin{equation}
\beta_{\rm burst}=\frac{\#\{i:\std(B_i)>1.5\,\med_k\std(B_k)\}}{\#\{i\}},\qquad
\text{CV}=\frac{\std_i\,\std(B_i)}{\operatorname{mean}_i\std(B_i)},
\end{equation}
and the drift is the slope of a straight line through $\ell_j$ against time (ADU/ms).

\paragraph{Clocking-induced noise.} The in-window residual spectrum $P_w(f)$ (Hann, window length $L_w-10$) is
compared with the spectrum $P_g(f)$ of the last 600 samples of the static gaps after H transfers (no clock switching):
\begin{equation}
\mathcal C=\sqrt{\frac{\overline{P_w}/(L_w-10)}{\overline{P_g}/600}}\Big|_{0.2\text{--}3\,\text{MHz}} ,
\end{equation}
which is 1 if switching adds no noise.

%=====================================================================
\section{Separating white and 1/f noise with the CDS window}\label{sec:H}

\subsection{The transfer function}
The per-skip statistic $\sigma_1^2$ is a linear filter of the input noise. For a stationary noise of one-sided PSD
$S(f)$:
\begin{equation}
\sigma_1^2=\int_0^\infty S(f)\,|H(f)|^2\dd f .
\end{equation}
Each factor of $H$ is one operation:
\begin{align}
\text{boxcar of length }\tau:&\quad \sinc^2(f\tau)=\Big(\tfrac{\sin\pi f\tau}{\pi f\tau}\Big)^2\\
\text{signal minus pedestal, centres } \Delta t \text{ apart:}&\quad |1-e^{-2\pi i f\Delta t}|^2=4\sin^2(\pi f\Delta t)\\
\text{consecutive-skip difference}/\sqrt2:&\quad \tfrac12|1-e^{-2\pi i fT}|^2=2\sin^2(\pi fT)
\end{align}
\begin{equation}
\boxed{|H(f)|^2=4\sin^2(\pi f\Delta t)\,\sinc^2(f\tau)\,2\sin^2(\pi fT)}
\end{equation}
Check: for white $S_w$, the $\sin^2$ factors average to $\tfrac12$ each and $\int_0^\infty\sinc^2(f\tau)\dd f=1/(2\tau)$,
so $\sigma_1^2=S_w/\tau$, i.e.\ eq.~\eqref{eq:en}. At low $f$, $|H|^2\propto f^4$, so the 1/f integral converges.

\subsection{The model and the fit}
With $S(f)=S_w+K/f$:
\begin{equation}
\sigma_1^2(\tau)=S_w\,I_w(\tau)+K\,I_f(\tau),\qquad
I_w=\int|H|^2\dd f,\quad I_f=\int\frac{|H|^2}{f}\dd f .
\end{equation}
The integrals are evaluated numerically on a log grid (1\,Hz--7.5\,MHz). White noise scales as $I_w\propto1/\tau$;
the 1/f term is nearly independent of $\tau$. Measuring $\sigma_1$ at several $\tau$ and solving the weighted linear
least-squares problem (weights $1/\sigma_1^2$, i.e.\ relative errors) gives $(S_w,K)$. Derived:
\begin{align}
e_w&=\sqrt{S_w}/G &&\text{white part (nV/}\sqrt{\text{Hz}}\text{): 12.6--14.2 PA03, 7.5--9.1 LSM}\\
f_c&=K/S_w &&\text{1/f corner (where } K/f=S_w\text{): 16--39\,kHz PA03, 61--334\,kHz LSM}\\
\phi_{1/f}(\tau)&=\frac{K I_f}{S_w I_w+K I_f}\Big|_{\tau=10\,\mu\text{s}} &&\text{1/f share of the single-skip variance}
\end{align}
and the fit quality $\sqrt{\langle(\hat\sigma_1^2/\sigma_1^2-1)^2\rangle}$.

\subsection{Three ways of varying the window}
\begin{center}\small
\begin{tabularx}{\linewidth}{@{}lXX@{}}\toprule
& how $\tau$ varies & $\Delta t$, $T$ in $|H|^2$\\\midrule
software (\code{trap\_from\_trace.py}) &
from one raw trace: the \emph{last} $\tau$ samples of each 10\,\textmu s window, $\tau=30,45,60,90,120,140$ samples
(2--9.3\,\textmu s); pedestal $\frac1\tau\sum_{m=L_w-\tau}^{L_w-1}c_j[m]$, signal likewise from $s_0$ &
$\Delta t=s_0/f_s=13.8$\,\textmu s fixed; $T=P/f_s$\\
hardware (\code{oxide\_traps\_tau.py}) & 100 v2 strips, \code{delayInt} 300--3000 cycles & $\Delta t=\tau+3.8$\,\textmu s;
$T=32.6+2(\tau-10)$\,\textmu s\\
skip averaging (\code{oxide\_traps.py}) & $\tau$ fixed, average $N$ skips & skip-difference factor replaced by the $N$-skip average, eq.~\eqref{eq:HN}\\\bottomrule
\end{tabularx}
\end{center}
\begin{equation}
|H_N(f)|^2=4\sin^2(\pi f\Delta t)\,\sinc^2(f\tau)\Big[\frac{\sin\pi fNT}{N\sin\pi fT}\Big]^2,\qquad
\sigma_N^2=S_wI_w(N)+K I_f(N).\label{eq:HN}
\end{equation}
The $N$-averaging variant is weakly constraining: through the CDS, 1/f noise averages down almost like white noise, so
$\sigma_N$ alone cannot separate them. The window variants are the ones used for the table.

\paragraph{Own vs total noise.} With simultaneous strips (hardware variant) each channel's per-skip series is first
cleaned of the part shared with the other channels of its board:
\begin{equation}
u_c=x_c-\beta_c\,\bar x_{\neg c},\qquad \bar x_{\neg c}=\frac1{n-1}\sum_{o\ne c}x_o,\qquad
\beta_c=\frac{\langle x_c,\bar x_{\neg c}\rangle}{\langle\bar x_{\neg c},\bar x_{\neg c}\rangle}
\end{equation}
(pixel means removed first), and the fit is done on $u_c$ (``own'') and on $x_c$ (``total''). The software variant
from single traces fits the total.
\caveat{The fit attributes to $K$ everything that is $\tau$-independent: transistor 1/f, setup wander, and
band-limited white noise (the $\approx$0.9\,\textmu s settling makes very short windows non-white). The own white part
(9--13\,nV/$\sqrt{\rm Hz}$) sitting below the frontend's own noise is not yet explained. Until the fit is run on an open
input, $K$, $f_c$, $\phi_{1/f}$ and $N_t$ are \textbf{upper limits} on the transistor's own 1/f. The window optimum
(5--8\,\textmu s) uses only the fitted $\sigma_1(\tau)$ curve and is unaffected.}

%=====================================================================
\section{Oxide trap density and single-trap steps}

\paragraph{Referring $K$ to the gate.} $K$ is in ADU$^2$ at the ADC. In V$^2$ at the gate (= sense node):
\begin{equation}
K_{V_g}=\frac{K}{(G\cdot10^6\ \text{ADU/V})^2\,A^2}.
\end{equation}

\paragraph{McWhorter number fluctuation.} Traps within a tunnelling distance $\lambda_t$ of the interface exchange
carriers with the channel; their superposition gives
\begin{equation}
S_{V_g}(f)=\frac{q^2k_BT\lambda_t N_t}{WL\,C^2 f}\quad\Rightarrow\quad
\boxed{N_t=\frac{K_{V_g}\,WL\,C^2}{q^2k_BT\lambda_t}}\ \ [\text{J}^{-1}\text{cm}^{-3}]\ \times q\ \to\ [\text{eV}^{-1}\text{cm}^{-3}],
\end{equation}
with $C=C_{\rm ox}=0.454$\,fF/\textmu m$^2$ (table) or the buried-channel $C_{\rm eff}=0.125$ (alternative),
$A=0.85$ (table) or 0.44 (withdrawn kTC route), $WL=345$\,\textmu m$^2$, $T=130$\,K (PA03). Units in the code: $WL$ in
cm$^2$, $C$ in F/cm$^2$, $\lambda_t$ in cm. Result: $N_t\le(1.1$--$2.4)\times10^{16}$\,eV$^{-1}$cm$^{-3}$ (good
thermal oxide: $10^{16}$--$10^{17}$).
\caveat{$N_t\propto K/(G^2A^2)\cdot C^2\cdot WL/(T\lambda_t)$: a 15\% error in $C_{\rm ox}$ is 30\% in $N_t$;
$\lambda_t$ is a convention. The LSM values use $T=180$\,K (checked: $N_t/K$ is $180/130$ lower than for PA03).}

\paragraph{Telegraph (RTS) steps} (\code{oxide\_traps.py}). Running means of $w=25$ skips,
$\Delta_k=\bar u_{k+w}-\bar u_k$; a step is a local maximum of $|\Delta_k|>6\,\std(u)\sqrt{2/w}$. Step size of one
trapped charge, for comparison:
\begin{equation}
\delta_{1q}=\frac{q}{C\,WL}\,A\,G\cdot10^6\ \ [\text{ADU}].
\end{equation}

%=====================================================================
\section{Conversion gain and sense-node capacitance}

\begin{align}
\text{CG}_{\rm in}&=\frac{g}{G}\ \ [\text{\textmu V}/e^-] &&\text{at the video input: 1.52--1.62 (PA03), 1.44--1.58 (LSM)}\\
C_{\rm sn}&=\frac{A\,q}{\text{CG}_{\rm in}} &&\text{84--89\,fF (PA03), 86--95\,fF (LSM)}
\end{align}
Worked example: $g=1.37$\,ADU/e $\Rightarrow$ CG$_{\rm in}=1.566$\,\textmu V/e,
$C_{\rm sn}=0.85\times1.602\times10^{-19}/1.566\times10^{-6}=86.9$\,fF.

\paragraph{Follower relations used to justify $A$} (\code{chain\_calibration.py}, REPORT \S6):
\begin{equation}
A=\frac{g_mR_L}{1+g_mR_L},\qquad C_{\rm sn}=C_{\rm fixed}+C_{gs}(1-A),\qquad C_{gs}\approx\tfrac23C_{\rm ox}WL\approx100\text{--}109\ \text{fF}.
\end{equation}
With $R_L=20$\,k$\Omega$ and the measured CG$_{\rm in}=Aq/C_{\rm sn}$, each assumed $C_{\rm fixed}$ fixes $A$. Device
estimates ($g_m\approx170$--750\,\textmu S) give $A\approx0.8$--0.9, hence $C_{\rm fixed}\approx60$--80\,fF (``most of
the node is wiring'').

\paragraph{kTC route (withdrawn).} Each reset freezes charge $\sqrt{k_BTC_{\rm sn}}$, so
$C_{\rm sn}=(q\sigma_{\rm kTC})^2/(k_BT)$ with $\sigma_{\rm kTC}$ in electrons from the per-skip level scatter
$\sigma_{\rm lvl}$ minus the electronics part. The electronics part was taken from board-102 ch0--2 as ``dead''
channels; they have transistors, so the subtraction (and the 45\,fF / $A=0.44$ result) is withdrawn.

\caveat{$C_{\rm sn}$ is linear in $A$ and in $1/G$. The spread between amplifiers ($\pm6\%$) is robust (same $A$, $G$
for all); the absolute value is model-dependent.}

%=====================================================================
\section{Settling and bandwidth}

\paragraph{Exponential settling.} After the reset pulse (end of the bit-4-low run $r_1$) and after each H transfer,
the template (or the static-gap segment) is fitted with
\begin{equation}
y(t)=a+b\,e^{-t/\tau_s},\qquad 0.05\le\tau_s\le500\ \text{\textmu s},
\end{equation}
on $t\in[r_1+3,\ r_1+403]$ samples (reset) or the first 400 samples of the gap. $\tau_s\approx R_{\rm out}C_{\rm line}$:
0.9--1.1\,\textmu s (PA03), 0.3--1.0 (LSM).

\paragraph{In-window slope.} A straight line through the template inside each window (trimmed by 5 samples):
$\partial\bar c/\partial t$ in ADU/\textmu s, divided by $g$ for e/\textmu s. Non-zero means the window is not settled
(LSM board 109 after the summing-well pulse).

\paragraph{CDS offset (summing-well feedthrough).} $\overline{s}-\overline{p}$ of the template, $\div g$ in electrons.

\paragraph{Rise time} (\code{src/trace\_features.py}). 10--90\% rise of the reset pulse rebuilt by equivalent-time
sampling (samples placed at their exact time relative to the edge, 0.1-sample bins).

%=====================================================================
\section{Clock-line to sense-node capacitances}\label{sec:cpl}

\paragraph{Why charge per volt is a capacitance.} A step $\Delta V_k$ on line $k$, coupled by $C_k$ to a node of total
capacitance $C_{\rm sn}$, moves the node by $\Delta V_{\rm sn}=C_k\Delta V_k/C_{\rm sn}$. Through the conversion gain
$q/C_{\rm sn}$ this is the same as $N=C_k\Delta V_k/q$ electrons. So
\begin{equation}
\boxed{C_k=q\,\frac{N}{\Delta V_k}=q\,\frac{1}{g}\Big|\frac{\partial y}{\partial V_k}\Big|}\qquad(y\ \text{in ADU}),
\end{equation}
independent of $C_{\rm sn}$, $A$ and $G$: only the image gain $g$ enters.

\paragraph{Regression on all rails at once} (board 101 voltage scan, sessions G+H: every rail's high and low varied
together in a Sobol design). For each amplifier and each response $y$,
\begin{equation}
y_i=\beta_0+\sum_k\beta_k\,(V_{k,i}-\tilde V_k)+\varepsilon_i,\qquad
\operatorname{Var}\hat\beta=(X^\top X)^{-1}\hat\sigma^2_\varepsilon ,
\end{equation}
with rails that moved in lockstep ($|r|>0.999$) merged. Then $C_k=q|\hat\beta_k|/g$ and significance $z_k=\hat\beta_k/
\operatorname{se}(\hat\beta_k)$. The responses $y$:
\begin{itemize}
\item \textbf{settled CDS offset} of an empty skip (\code{sense\_node\_coupling.py}): only lines whose level differs
between the two windows (SW, OG, RG through the reset) contribute; video-cable pickup has decayed ($\sim$1\,\textmu s)
by the window, so this is pickup-free. SW: 0.08--0.13\,fF.
\item \textbf{edge jumps} (\code{clock\_edges.py}, \code{v\_edges.py}): the median waveform $w$ around the bit-6 (H) or
bit-7 (V) flags, edges $e$ located from the largest $|\Delta^2w|$, and
\begin{equation}
J_e=\frac19\sum_{m=3}^{11}w[e+m]-\frac19\sum_{m=-12}^{-4}w[e+m]
\end{equation}
(the fast, capacitive part before relaxation). H$_{1,2,3}$: 0.011--0.016\,fF ($z>4$); V$_{1,2,3}$: not detected,
$<0.003$\,fF (3$\sigma$); RG low rail at the reset edges: 0.15--0.22\,fF.
\item \textbf{feedthrough per volt} (relative, all setups): reset-pulse height above the signal level / $(\text{RG}_H-\text{RG}_L)$,
summing-well dip / SW swing, H/V waveform peak-to-peak / swing, each $\div g$ (e/V). These include board pickup and the
whole transfer waveform, so they compare channels, not absolute $C_k$.
\end{itemize}

\paragraph{Reset-gate overlap.}
\begin{equation}
L_{\rm ov}=\frac{C_{\rm RG}}{C_{\rm ox}W_{\rm RG}}=\frac{0.21\text{--}0.37\ \text{fF}}{0.454\ \text{fF/\textmu m}^2\times6\ \text{\textmu m}}
=0.08\text{--}0.14\ \text{\textmu m}\quad(\text{101 ch2: }1.7\times).
\end{equation}
\caveat{$L_{\rm ov}\propto t_{\rm eq}$ through $C_{\rm ox}$. The 0.21--0.37\,fF RG values from the feedthrough feature
had the board-102 pickup subtracted assuming those channels were dead; the trace-edge regression (0.15--0.22) does not
need that subtraction.}

%=====================================================================
\section{Reset charge injection and the buried-channel depth}

\paragraph{Measurement.} In the rail regression of the reset-edge jumps (and their relaxation), the coefficients of
RG$_H$ and VR come out equal and opposite (+17\,500 to +20\,200 and $-$16\,000 to $-$18\,700 ADU/V): the response depends
on $V_G-V_R$. With Sec.~\ref{sec:cpl},
\begin{equation}
C_{\rm inj}=\frac{q}{g}\,\frac{\partial J}{\partial(\text{RG}_H-\text{VR})}=2.1\text{--}2.4\ \text{fF/V}.
\end{equation}

\paragraph{Channel charge of the reset FET.} When the reset FET turns off, its inversion charge
\begin{equation}
Q_{\rm ch}=C_{\rm eff}\,W_{\rm RG}L_{\rm RG}\,(V_G-V_{\rm ch}-V_T)
\end{equation}
splits between source (sense node) and drain; with a fraction $\alpha\approx\tfrac12$ onto the node,
$C_{\rm inj}=\alpha C_{\rm eff}W_{\rm RG}L_{\rm RG}$, so
\begin{equation}
C_{\rm eff}=\frac{C_{\rm inj}}{\alpha\,W_{\rm RG}L_{\rm RG}}=\frac{2.1\text{--}2.4}{0.5\times36}=0.117\text{--}0.133\ \text{fF/\textmu m}^2 .
\end{equation}

\paragraph{Depth.} In a buried channel the gate couples through the dielectric and the silicon above the potential
maximum in series:
\begin{equation}
\frac1{C_{\rm eff}}=\frac{t_{\rm eq}}{\varepsilon_{\rm ox}}+\frac{x_{\rm ch}}{\varepsilon_{\rm Si}}
\ \Rightarrow\
x_{\rm ch}=\varepsilon_{\rm Si}\Big(\frac1{C_{\rm eff}}-\frac{t_{\rm eq}}{\varepsilon_{\rm ox}}\Big).
\end{equation}
Numbers ($\varepsilon_{\rm ox}=0.0345$, $\varepsilon_{\rm Si}=0.1036$\,fF/\textmu m): $1/C_{\rm eff}\approx8.0$\,\textmu
m$^2$/fF, $t_{\rm eq}/\varepsilon_{\rm ox}=0.076/0.0345=2.20$, remainder $5.8\times0.1036\Rightarrow
\boxed{x_{\rm ch}\approx0.6\ \text{\textmu m}}$ (0.3--0.9 for $\alpha$ and gate area $\pm30\%$; recipe: B 150\,keV,
projected range 0.45\,\textmu m plus diffusion).
\caveat{Only the combination $1/C_{\rm eff}$ is measured. $t_{\rm eq}$ and $x_{\rm ch}$ are degenerate: assuming
$x_{\rm ch}=0.45$\,\textmu m instead gives $t_{\rm eq}\approx126$\,nm. An LCR measurement of a clock-line capacitance
($C_{\rm line}=\varepsilon_{\rm ox}A/t_{\rm eq}$) breaks the degeneracy; an RG-edge-speed scan pins $\alpha$.}

%=====================================================================
\section{Knees and channel potentials (\code{channel\_knees.py})}

\paragraph{Hinge fit.} For each rail $x$ and output $y$ (log of $\sigma_e$, $\lambda$, $g$, CIC, white noise), with all
other rails as linear regressors $X$:
\begin{equation}
\ln y=X\beta+c\,\max(0,x-k)+\varepsilon,
\end{equation}
$k$ by grid search (33 quantiles of $x$, 10--90\%) minimising the residual sum of squares $\text{RSS}_1(k)$. Against the
straight line ($\text{RSS}_0$):
\begin{equation}
F=\frac{\text{RSS}_0-\text{RSS}_1}{\text{RSS}_1/(n-p-2)}\sim F_{2,\,n-p-2},
\end{equation}
kept if $p<0.01$ and found on $\ge2$ CCDs; error on $k$ by bootstrap over packets (60 resamples).

\paragraph{Dose shift from a knee offset.} A change $\Delta N_s$ of the buried-channel sheet dose shifts the channel
potential by
\begin{equation}
\Delta V\approx q\,\Delta N_s\Big(\frac{t_{\rm eq}}{\varepsilon_{\rm ox}}+\frac{x_j}{\varepsilon_{\rm Si}}\Big)
=\Delta N_s\times1.05\times10^{-11}\ \text{V\,cm}^2\quad(t_{\rm eq}=76\ \text{nm},\ x_j=0.45\ \text{\textmu m}).
\end{equation}
CCD\_B's H3-low knee sits $0.85\pm0.3$\,V below CCD\_C/D $\Rightarrow\Delta N_s\approx8\times10^{10}$\,cm$^{-2}$
$\approx0.6\%$ of $1.3\times10^{13}$ (or the same density of fixed oxide charge).

%=====================================================================
\section{The setup: shared noise, crosstalk}

\paragraph{Shared vs own} (\code{common\_mode.py}, \code{setup\_noise\_index.py}). Per channel, per-skip CDS with pixel
means removed (optionally the pixel-averaged droop pattern too), averaged in blocks of $n$ skips ($y$), and the
projection on the mean of the board's other channels ($\bar y_{\neg c}$) as in Sec.~\ref{sec:H}:
\begin{align}
\text{shared share}&=1-\frac{\operatorname{Var}(y-\beta\bar y_{\neg c})}{\operatorname{Var}y},\\
\text{SNI}_n&=\sqrt{\frac{\operatorname{Var}y-\operatorname{Var}(y-\beta\bar y_{\neg c})}{\operatorname{Var}(y-\beta\bar y_{\neg c})}},\qquad
X_n=\sqrt{1+\text{SNI}_n^2}=\frac{\sigma_{\rm total}}{\sigma_{\rm own}} .
\end{align}
SNI = 0: no setup noise; 1: the setup adds as much as the amplifier. Removing the shared part entirely would lower the
per-skip noise by $1-1/X_1$ (v3: 9.46$\to$6.5\,ADU, $-30\%$).
\caveat{``Shared'' is defined against the other channels of the same board. Noise that is common to all four (e.g.\ a
shared return) is captured; noise entering one channel alone (its own cable, its own frontend) counts as ``own''. Applied to the per-channel \emph{traces}, which are separate
acquisitions, the shared part is a sequence-locked pattern, not simultaneous noise (Sec.~\ref{sec:setup2}); the
trace-based SNI values are withdrawn as setup numbers.}

\paragraph{Crosstalk} (\code{crosstalk.py}, images). Source pixels with $q_A>30\,e^-$, target pixels without own
charge ($|q_B|<4\sigma_B$, $\sigma_B$ from the MAD); $q_B$ vs $q_A$ binned in 20 quantiles, Theil--Sen slope of the bin
medians:
\begin{equation}
k_{A\to B}=\frac{\partial q_B}{\partial q_A}\quad(\sim10^{-4}).
\end{equation}

\paragraph{Mutual capacitance from the window dependence} (\code{mutual\_capacitance.py}). With crosstalk at several
$\tau$,
\begin{equation}
k(\tau)=\frac a\tau+b,\qquad C_m^{\rm video}=\frac a{R_{\rm out}}\ (R_{\rm out}=2.5\text{--}5\ \text{k}\Omega),\qquad
C_m^{\rm sn}=b\,C_{\rm sn} :
\end{equation}
a step on one video line injects a pulse of area $R_{\rm out}C_m\Delta V$ into the other, which a window of length
$\tau$ averages to $R_{\rm out}C_m\Delta V/\tau$; the $\tau$-independent part is coupling before the followers (or a
shared impedance, if negative).

%=====================================================================
\section{What characterises the setup, and the cabling change}\label{sec:setup2}

By \emph{the setup} we mean everything between the output transistor and the number written to disk: the video line
on the flex and cable, the frontend board (20\,k$\Omega$ load, preamplifier, AC coupling), the ADC and its clock, and
the grounds, supplies and bias lines the channels share. A number characterises the setup only if the amplifier's own
physics cannot produce it. This section says how each setup number is isolated from the amplifier, why the familiar
per-skip noise $\sigma_1$ is \emph{not} one of them, and works through the one setup change in these data: the
grounding/cabling change made between the October runs (Oct.\ 6--12) and the November runs (from Nov.\ 2).

%---------------------------------------------------------------------
\subsection{How a number is separated from the amplifier}

\begin{center}\small
\begin{tabularx}{\linewidth}{@{}lXX@{}}\toprule
handle & why the amplifier cannot produce it & used for\\\midrule
before/after, same CCDs and voltages &
the transistors, their traps and their sense nodes do not change when a cable is moved; anything that moves at fixed
voltages is the setup. Amplifier-only numbers are \emph{controls} and must stay put &
the cabling change (Sec.~\ref{sec:cabling})\\
simultaneity &
four output transistors on four CCDs have independent white and 1/f noise; a part that is correlated \emph{in the same
instant} on all of them must enter through something they share (return, supply, bias, clock lines, cable) &
$\rho_{\rm sim}$, common-mode gain, strip shared/own ratio\\
same on a board, different between boards &
a circuit constant of the readout path, not of a transistor (the four amplifiers of a board are on different CCDs) &
ringing, baseline recovery, gain-scale step\\
independent of the CCD voltages &
a transistor-limited quantity follows its operating point (VDD, VR); one that does not is set by the interconnect &
rise time ($-0.1$ to $+0.8$\,\%/V in VDD)\\
spectral shape &
transistor noise is featureless (white + 1/f); narrow lines are interference &
the 2.5\,kHz line and its harmonic\\\bottomrule
\end{tabularx}
\end{center}
\caveat{The per-channel \emph{traces} are separate acquisitions about 1\,s apart (the other three streams of each file
are empty frames). Their cross-channel ``shared'' part is a pattern locked to the readout sequence, not simultaneous
noise, so the simultaneity handle needs images or level-1 strips, which record the four amplifiers together. A trace on
its own can only use the other handles.}

%---------------------------------------------------------------------
\subsection{$\sigma_1$ in detail, and why it does not see the setup}\label{sec:sigma1}

\paragraph{What $\sigma_1$ is.} A skipper amplifier measures the same pixel charge $Q$ many times without destroying it.
Skip $j$ of a pixel gives $x_j=gQ+n_j$, where $n_j$ is the noise of that one CDS measurement (pedestal window minus
signal window). Because $Q$ is identical in every skip of the pixel, the difference of two consecutive skips contains
\emph{no charge}:
\begin{equation}
d_j=x_{j+1}-x_j=n_{j+1}-n_j,\qquad
\operatorname{Var}d_j=2\sigma_n^2\ \text{(if consecutive } n_j \text{ are independent)}
\ \Rightarrow\ \sigma_1=\frac{\std(d_j)}{\sqrt2}.
\end{equation}
So $\sigma_1$ is the noise of a single CDS measurement, obtained from a raw trace with no image, no gain fit and no
knowledge of the charge (differences across a pixel boundary are excluded, since $Q$ changes there). It is in ADU;
$\sigma_1/g$ is electrons and $\sigma_1\sqrt\tau/G$ an input density (eq.~\eqref{eq:en}).

\paragraph{The textbook route.} Usually the CDS noise is \emph{predicted} from a noise spectrum $S(f)$ taken with the
clocks stopped (or from a device model), by integrating it through the CDS filter:
\begin{equation}
\sigma_{\rm CDS}^2=\int_0^\infty S(f)\,|H_{\rm CDS}(f)|^2\dd f,\qquad |H_{\rm CDS}|^2=4\sin^2(\pi f\Delta t)\,\sinc^2(f\tau).
\end{equation}

\paragraph{How ours differs.}
\begin{enumerate}
\item \textbf{Measured, not integrated.} $\sigma_1$ is read off the clocked trace in the time domain. It contains every
noise the clocking adds (clock-induced noise, bursts, sampling jitter) and needs no stationarity. The integral is used
only to \emph{interpret} it (Sec.~\ref{sec:H}).
\item \textbf{One extra filter factor.} The skip difference multiplies the filter by $2\sin^2(\pi fT)$ (Sec.~\ref{sec:H}).
For white noise this averages to 1, so $\sigma_1=\sigma_{\rm CDS}$. At low frequency it does not:
$|H_{\rm CDS}|^2\approx4(\pi f\Delta t)^2\propto f^2$, while $|H|^2\propto f^4$. At 200\,Hz the extra factor is
$8.4\times10^{-4}$, at 1\,kHz $2.1\times10^{-2}$.
\end{enumerate}

\paragraph{Consequence: $\sigma_1$ is an amplifier number.} The share of each statistic's variance that a flat noise,
at the same density as the white floor, contributes from a band,
$w_{\rm band}=\int_{\rm band}|H|^2\dd f\big/\int_0^\infty|H|^2\dd f$ (computed with $\Delta t=13.8$, $\tau=10$,
$T=32.6$\,\textmu s):
\begin{center}\small
\begin{tabular}{lccc}\toprule
band & $\sigma_1$ (skip difference) & one CDS & 500-skip pixel, $|H_{500}|^2$ of eq.~\eqref{eq:HN}\\\midrule
20--200\,Hz & $1.0\times10^{-10}$ & $2.0\times10^{-7}$ & $1.3\times10^{-6}$\\
200\,Hz--1\,kHz & $3.2\times10^{-7}$ & $2.5\times10^{-5}$ & $5.8\times10^{-6}$\\
1--10\,kHz & $2.3\times10^{-2}$ & $2.4\times10^{-2}$ & $7.1\times10^{-5}$\\\bottomrule
\end{tabular}
\end{center}
Relative to its own white noise, a 500-skip pixel is $\sim10^4$ times more sensitive to 20--200\,Hz noise than
$\sigma_1$. Setup noise (ground loops, supply and bias wander) lives at these low frequencies, so it shows up in
$N$-skip images and long-$N$ strip statistics, and hardly at all in $\sigma_1$. This is why $\sigma_1$ serves as a
\emph{control} for a setup change: across the cabling change it moved by $-3\,\%$ while the image noise fell by 30\,\%.
\caveat{The 500-skip column uses one pixel's filter, without the pixel-to-pixel difference and row-baseline subtraction
that the image analysis applies (these remove even slower noise). The orders of magnitude, not the exact values, are
the point.}

%---------------------------------------------------------------------
\subsection{The ASD and why the trace one is only a relative setup number}

\paragraph{Definition.} The power spectral density $S_{rr}(f)$ (one-sided, ADU$^2$/Hz, so $\int_0^\infty S_{rr}\dd f$
is the variance of $r$) is estimated by Welch's method; the amplitude spectral density is its square root, referred to
the video input with the chain gain:
\begin{equation}
\text{ASD}(f)=\frac{\sqrt{S_{rr}(f)}}{G}\qquad[\text{V}/\sqrt{\text{Hz}}].
\end{equation}
It says how much noise amplitude sits at each frequency. A narrow line is interference; the 200\,Hz--1\,kHz level is
where setup wander sits.

\paragraph{Why it is not integrated to a $\sigma$.} On a clocked trace, $r$ is the cycle minus the median template.
The ADC is not phase-locked to the sequencer (Sec.~\ref{sec:circuit}), so every reset pulse ($\sim$50\,kADU) is
sampled at a random phase and leaves a large residual, and the cycle joins add steps. Integrating that spectrum through
the CDS filter gives $\sigma_1\approx1400$\,e on PA03 (\code{setup\_noise\_N.py}), which is meaningless. The clocked
ASD is therefore used only \emph{like for like}: the same sequencer and the same trace length, before and after a
change. The textbook forward integration works on \emph{static} records (all clocks frozen, LSM): there
$\sigma^2_{N,\rm band}=\int_{\rm band}S(f)|H_N(f)|^2\dd f$ predicts the measured per-skip $\sigma$ to a median ratio of
0.97, and gives the setup's share of a 500-skip pixel in electrons (0.012--0.14\,e on the LSM channels).

%---------------------------------------------------------------------
\subsection{The setup numbers, as computed}\label{sec:circuit}

\paragraph{Simultaneous correlation and common-mode gain} (\code{oct\_nov.py}, images; the four amplifiers of a board
in one file). Per image, each amplifier's map minus its median, then minus each row's median (``row baseline''). Robust
width $s_c=1.4826\,\med|X_c|$. Only pixels where \emph{every} amplifier is in its 0-e$^-$ peak ($|X_c|<2.5s_c$) are
kept, so no charge enters. Then
\begin{align}
\rho_{\rm sim,c}&=\tfrac13\sum_{o\ne c}\operatorname{corr}(X_c,X_o),\\
\text{CMS gain}_c&=1-\frac{\sigma\big(X_c-\beta_c\bar X_{\neg c}\big)}{\sigma(X_c)},\quad
\beta_c=\frac{\langle X_c,\bar X_{\neg c}\rangle}{\langle\bar X_{\neg c},\bar X_{\neg c}\rangle}
\ \text{fitted on alternate pixels, applied to the others,}
\end{align}
with $\sigma$ the MAD width and $\sigma_{0e,c}=s_c/g_c$ in electrons.
\emph{Splitting a pixel's noise.} If each amplifier sees $X_c=s+n_c$, with a common part $s$ and independent own
noise $n_c$ of similar size on every channel, then $\rho_{\rm sim}=\operatorname{Var}s/\operatorname{Var}X$, so
\begin{equation}
\sigma_{\rm common}=\sqrt{\rho_{\rm sim}}\;\sigma_{0e},\qquad \sigma_{\rm own}=\sqrt{1-\rho_{\rm sim}}\;\sigma_{0e}.
\end{equation}

\paragraph{Shared/own power ratio from a level-1 strip.} A strip records all four amplifiers together and keeps every
skip. Per channel the per-skip CDS minus its pixel mean is $r_c[j]$; the shared series is the channel mean
$\bar r[j]=\frac14\sum_c r_c[j]$ and the own series $r_c-\bar r$. Welch spectra of both give
\begin{equation}
R_{\rm band}=\frac{\int_{\rm band}S_{\bar r\bar r}\dd f}{\big\langle\int_{\rm band}S_{(r_c-\bar r)(r_c-\bar r)}\dd f\big\rangle_c},
\end{equation}
reported for 0--200\,Hz, 0.2--1\,kHz and 1--5\,kHz, together with the narrow lines present on every channel.
\caveat{The band ratios were computed in an analysis recorded in REPORT \S24; \code{strip\_noise.py} computes the
time-domain counterpart (correlation $\rho_N$ and shared fraction of $N$-skip block differences).}

\paragraph{Readout-chain scale.} With the image gains of two periods,
$r_c=g_c^{\rm after}/g_c^{\rm before}$. A step equal on all four amplifiers (four CCDs, four transistors) is in the
chain they share, not in the sense nodes.

\paragraph{Circuit constants of the readout path} (\code{trace\_physics2.py}, \code{src/trace\_features.py}).
\begin{itemize}
\item \textbf{Baseline recovery.} The median, over pixels, of the pedestal level against time within a pixel, fitted
with $\ell(t)=a+b\,e^{-t/\tau_{\rm rec}}$: the video chain's AC coupling recharging after each H transfer.
Board 101: $\tau_{\rm rec}=8.2\pm0.3$\,ms, $b=-1.8$\,kADU, identical on all four channels; board 102: none
($\tau\ge130$\,ms). The within-pixel CDS drift has the same 8\,ms constant.
\item \textbf{Baseline leakage into the CDS.} The per-skip CDS regressed on the per-skip level within a pixel gives the
fraction $\varepsilon$ of a level fluctuation that survives the CDS; rejection $=20\log_{10}(1/|\varepsilon|)$:
$\varepsilon=-0.18\,\%$ (101, 55\,dB), $+0.48\,\%$ (102, 46\,dB).
\item \textbf{Sampling-phase jitter.} The half-height crossing time of each reset pulse (linear interpolation), its
consecutive differences $/\sqrt2$: 20--25\,ns. A sampling clock that is not locked to the sequencer gives a uniformly
random phase, $1/(f_s\sqrt{12})=19.2$\,ns. So the ADC runs asynchronously to the 10\,ns sequencer clock: a firmware
property of the setup.
\item \textbf{Ringing.} The strongest frequency in the residual of the reset-edge fit, over the 25 samples after the
peak (zero-padded spectrum, peak refined by a parabola): 3.1\,MHz on board 101 (160--240\,e rms), 0.55--0.7\,MHz on
board 102. It is the same on a board's four amplifiers, differs between boards, and the rise time does not follow VDD,
so it is an LC resonance of the readout path, not the transistor.
\end{itemize}
\caveat{25 samples span 1.67\,\textmu s, less than one period at 0.55--0.7\,MHz: board 102's ringing is ``much slower
than board 101'', not a precise frequency.}

%---------------------------------------------------------------------
\subsection{Worked example: the cabling change}\label{sec:cabling}

\paragraph{The data.} Board 101, 500-skip unbinned images at the October nominal voltages (VSUB 45.1\,V, VR
$-3.0$\,V, VDD $-19.0$\,V, H1 high 9.5\,V, V1 high 9.9\,V, H delay 1800): 12 October images, 19 from Nov.\ 10--12
(medians below; $\rho_{\rm sim}$ and CMS gain with the row baseline removed). Traces of equal length (0.2\,s, three in October). One October level-1 strip
($100\times1000\times3000$) and the November v3 strips.

\begin{center}\small
\begin{tabularx}{\linewidth}{@{}>{\raggedright\arraybackslash}Xccl@{}}\toprule
quantity & October & November & role\\\midrule
\multicolumn{4}{@{}l}{\emph{setup numbers}}\\
$\rho_{\rm sim}$, ch0--3 & 0.28 / 0.26 / 0.19 / 0.22 & 0.07 / 0.07 / 0.05 / 0.02 & simultaneity\\
CMS gain, ch0--3 & 8.0 / 6.3 / 4.8 / 6.2\,\% & 0.1 / 0.0 / 0.2 / 0.2\,\% & simultaneity\\
strip shared/own, 0--200\,Hz & 23 & 3.9 & simultaneity\\
strip shared/own, 0.2--1 / 1--5\,kHz & 1.3 / 1.6 & 0.7 / 1.5 & simultaneity\\
trace ASD 0.2--1\,kHz & 48.8\,\textmu V/$\sqrt{\rm Hz}$ & 34.8 ($-29\,\%$) & before/after\\
trace ASD 1--10 / 10--100\,kHz & 5.2 / ref. & 3.7 ($-28\,\%$) / $-15\,\%$ & before/after\\
gain scale $r_c$, ch0--3 & 1 (ref.) & 1.056 / 1.061 / 1.067 / 1.064 & same on a board\\
shared narrow lines & 2532, 5064\,Hz & 2240--2490, 5050--5240\,Hz & spectral shape\\\midrule
\multicolumn{4}{@{}l}{\emph{result}}\\
image $\sigma_{0e}$, ch0 / ch1 & 0.207 / 0.210\,e & 0.145 / 0.148\,e ($-30\,\%$) & \\\midrule
\multicolumn{4}{@{}l}{\emph{controls (amplifier numbers)}}\\
$\sigma_1$ (per skip) & ref. & $-3\,\%$ & \\
white density ch0/1/3 (nV/$\sqrt{\rm Hz}$) & 20.2--20.7 / 20.8--21.2 / 23.7--23.8 & 19.8 / 20.2 / 21.9 & \\\bottomrule
\end{tabularx}
\end{center}

\paragraph{Step 1: the amplifiers did not change.} $\sigma_1$ moved by $-3\,\%$ and the white density by 2--8\,\%.
Both are dominated by the transistor and the frontend's white noise (Sec.~\ref{sec:sigma1}), so the amplifiers are the
same before and after.

\paragraph{Step 2: the pixel noise did change.} At the same voltages, ch0 went from 0.207 to 0.145\,e
($\sigma^2$: 0.0428 $\to$ 0.0210\,e$^2$, half the variance gone).

\paragraph{Step 3: how much of it was shared.} With the split above, ch0 in October:
$\sigma_{\rm common}=\sqrt{0.279}\times0.207=0.109$\,e and $\sigma_{\rm own}=\sqrt{0.721}\times0.207=0.176$\,e.
November: $\sqrt{0.070}\times0.145=0.038$\,e and $\sqrt{0.930}\times0.145=0.140$\,e. The common variance fell
$\times8.1$ (ch1: 0.107 $\to$ 0.039\,e, $\times7.4$). The strips say the same thing in frequency: below 200\,Hz the
shared power fell from 23 to 3.9 times the own power, about $\times6$; above 1\,kHz it did not change.

\paragraph{Step 4: the ``own'' part fell too.} $\sigma_{\rm own}$ went from 0.176 to 0.140\,e ($-20\,\%$; ch1 $-21\,\%$)
while the amplifier controls moved by only a few per cent. A correlation counts noise picked up by one channel's own
cable as ``own''; the before/after handle shows that this part was also setup. The change therefore removed both the
common slow noise (eight times less variance) and some per-channel pickup.

\paragraph{Step 5: what is left.} In November the common part is 0.04\,e and subtracting it gains nothing (CMS gain
$\approx0$). The 2.5\,kHz line and its harmonic are on every channel in both months, so the cabling did not touch
their source (a switching converter or a periodic digital process). The 500-skip pixel (16.3\,ms) averages them out;
other skip counts may not.

\paragraph{The recommended setup numbers.} From one image: $\rho_{\rm sim}$ (target 0) and the CMS gain. From one
$\sim$30\,s strip: the shared/own ratio below 200\,Hz (target 0) and the list of shared lines. Always quote $\sigma_1$
and the white density alongside, as the controls showing the amplifiers did not change.
\caveat{(i) The October image gains are 6\,\% lower ($r_c$); $\sigma_{0e}$ uses each period's gain, so the 30\,\% is in
electrons. The step is equal on all four channels, consistent with a chain change made at the same time; the trace
densities use the design $G$ in both months, so the 2--8\,\% control figure carries that uncertainty. (ii) The split
assumes the common part couples equally into all four channels. (iii) The October strip is one acquisition with
3000-skip pixels (November 480), so 23 vs 3.9 is indicative. (iv) October images had 10--100\,s exposures (November
0--1\,s); at 40--75\,e/pix/day this is $\le0.1$\,e/pix and the 0-e$^-$ selection excludes it. (v) ch2 (own telegraph
problem) and ch3 also improved but scatter more and are not used in the example.}

%=====================================================================
\section{Summary: which inputs each number depends on}

\begin{center}\small
\begin{tabularx}{\linewidth}{@{}lXl@{}}\toprule
quantity & depends on & status\\\midrule
$\sigma_1$ (ADU, e), $\sigma_N$, $\kappa$, $\rho_1$, $\mathcal R_{\rm img}$ & trace (+ $g$) & measured\\
$e_n$, ASD, setup number & $G$ (design) & measured, $G$-scaled\\
$f_c$, $\phi_{1/f}$ & window fit & upper limit (open-input test pending)\\
$e_w$ & window fit, $G$ & model\\
$N_t$ & $K$, $G^2A^2$, $C^2$, $WL$, $T$, $\lambda_t$ & upper limit, model\\
CG$_{\rm in}$ & $g$, $G$ & measured, $G$-scaled\\
$C_{\rm sn}$ & $g$, $G$, $A$ & model ($\pm6\%$ spread robust)\\
$C_k$ (RG, SW, H, V) & regression, $g$ & measured\\
$L_{\rm ov}$ & $C_{\rm RG}$, $C_{\rm ox}$, $W_{\rm RG}$ & model ($\propto t_{\rm eq}$)\\
$C_{\rm eff}$, $x_{\rm ch}$ & $C_{\rm inj}$, $\alpha$, $W_{\rm RG}L_{\rm RG}$, $t_{\rm eq}$ & model ($\pm30\%$, degenerate with $t_{\rm eq}$)\\
knee $k$, $\Delta N_s$ & hinge fit; $t_{\rm eq}$, $x_j$ & knee measured; dose model\\
$\tau_s$, slopes, CDS offset & template & measured\\
SNI, shared share, $k_{A\to B}$ & simultaneous channels / images & measured\\
$\rho_{\rm sim}$, CMS gain, strip shared/own ratio & simultaneous images / strips & measured (setup)\\
clocked-trace ASD bands & trace, $G$ & relative only (equal length, same sequencer)\\
$\tau_{\rm rec}$, leakage $\varepsilon$, phase jitter, ringing & trace & measured (setup constants)\\
$C_{\rm sn}$ from kTC & ``dead'' board-102 channels & withdrawn\\\bottomrule
\end{tabularx}
\end{center}

\end{document}
